\chapter*{附录四\quad 释文语译}
\phantomsection
\addcontentsline{toc}{chapter}{附录四\quad 释文语译}
\markboth{附录四}{附录四}

\noindent
下列三十七条，为释文之白话译。译者所加，非辑校者原文。

凡例三条：一、译文只求读通，不求文采。二、原件没有说的，译文不补；凡释文残缺处，译作「缺」。三、译文与释文有出入时，以释文为准。

\begingroup
\footnotesize
\linespread{1.32}\selectfont
\newcommand{\yi}[2]{\par\noindent{\sffamily\bfseries #1}\quad #2\par\vspace{2.2mm}}

\yi{简一}{一座山的下面，离海三舍远。地上没有草木，风从西边来，一年到头不停。用石头砌墙，墙厚四丈，高十二丈。里面是环，每环有一道门，各门不能同时打开。环外是路，路宽十步，绕环一圈。}

\yi{简二}{三环以上住的人少。一道门里，平常不超过一千人。环的尽头是界，界外不宜前往。界内每一步都在名册上。}

\yi{简三}{一环。二环。三环。四环。五环。六环。每环各一道门，门各有名册。名册由门吏保管，不入总库。}

\yi{简四}{气从下往上走，水从上往下走，都走管道。管道砌在墙里，看不见。管道每年检查一次，检查的人不进环，在环外的路上做事。查完把数据记在板上，板挂在门口。二十三年，三环西边路上的管道裂了。检查的人先发现，那一环没有关。赏了他，因为他没有把这事说出去。}

\yi{简五}{外工三百二十人，住在墙外的小屋。每月进环一次，进去不超过三个时辰。出来要洗澡，洗澡水不倒进渠里。衣服另外存放，不与环内的人混在一起。有一名外工越界，被抓住。问他为什么，他说：想看看山。按律处罚，名册没有除名。}

\yi{简六}{墙周长九里，高十二丈。墙外没有围墙，以地面为界。界线以外一百步，不许放东西，不许盖房，不许安葬。界线以内曾经有三处私盖的房子，拆了，盖房的人改了名册。}

\yi{简七}{第几环的西边储物间，有架子十二个，每个五层。层用板隔开，板下有小牌。进这个屋子的人要两个一起，一个记录，一个取物。}

\yi{简八}{凡东西入库，先记在板上，后记在册上。板与册不一样时，以板为准。板坏了就把字挪走，不挪板。册坏了就重抄，重抄的人署名。}

\yi{简九}{界这种东西，不是墙，不是壁，不是水，不是火。在里面的人出去，必定失去一半。失去的回不来。界外有山，山的名字不能叫。叫的人和越界同罪。}

\yi{简十}{每年统计人口，共三十二万一千。匠六万，工十一万，其余是居民。每三口人出一个工。工不入匠册，匠不入工册。}

\yi{简十一}{匠的名册单独一本。工的名册单独一本。居民不入册，只按人头计。三本册子各有各的吏，彼此不通气。}

\yi{简十二}{粟够九年，水够六年，气够三年。三样同时算，但不是同时用完。曾经有一年，粟只够三个月。削减从上面开始，居民每天一升，工照旧。}

\yi{简十三}{新来的人，先在外屋住三个月，才能进环。进环那天，给名册，给衣服，给器物。器物不够就用旧的补。旧器上有记号，记着上一个主人是谁。}

\yi{简十四}{额定人口三十一万二千，实际居住十一万九千。粟够九年，水够六年，气够三年。}

\yi{简十五}{每月发粟，每人每天二升。匠加倍。外工再加倍。不够就减，从上面开始减。}

\yi{简十六}{工的分配：捣的三千，搬运的一千二百，修理的八百，其余是杂役。捣的人不进环，口粮另行发放。}

\yi{简十七}{定额：每人每天三合。多的赏，少的罚。罚不超过三次，超过三次就改属别的名册。赏用粟，罚用劳役。十四年，捣的人里有十一个连着几天不够数，罚劳役三天，不改册。其中一个又犯，改了册。}

\yi{简十八}{名册每年换一次。换册那天，旧册烧掉。烧的两人，看的三人，都要签名。烧剩的灰埋在墙外。十九年，有一册没有烧。烧的人自己报了，看的三人说法一致。罚了烧的人，那一册最终没有烧。}

\yi{简十九}{事故的数目：兔乱一万一千，其他事故八千，逃亡三千，共二万二千。逃亡的人，从册上除名，名字削掉。}

\yi{简二十}{元年，清点人口。在册十一万九千，可以走的十一万九千，留下的没有。凡是不走的，名字不进这本册子。}

\yi{简二十一}{退役清单，共十九项。第一项：六处居住环，封闭。第二项：气路，切断。第三项：水路，切断。第四项：仓库，封闭。第五项：名册，烧毁。第六项：器物，毁掉。第七项：垂直运输装置（电梯）已拆，支柱在四十米高处截断，截断处做了记号，不要挖。第八项：界，切断。第九项：外屋，推平。第十项：墙门，堵死。}

\yi{简二十二}{第十一项：山，封闭。第十二项：炉子，熄火。第十三项：药，丢弃。第十四项：册子，收走。第十五项：界石，移走。第十六项：标记，抹去。第十七项：名字，削掉。第十八项：人口数，削减。第十九项：门，关闭。}

\yi{简二十三}{支柱截断，深四十米。填土，夯实。上面盖石板三层。石板上种草，草长不活。草长不活的有三处，其中一处在现在墙的西边。留守的人说：这三处，冬天不结冰。}

\yi{简二十四}{装船：人十一万九千。粟，全部带走。水，全部带走。气，全部带走。器物，挑要紧的三千件。册子，全部带走。凡是写着「全部」的，一件不留。}

\yi{简二十五}{凡是要走的人，各发一枚符。符不交回。符上刻他的名字。没有名字的不许上船。符丢了的不补。}

\yi{简二十六}{世代存活率：第一代六一四，第二代三八七，第三代一二九，第四代四二。四代之后没有记录。没有记录，不等于没有。}

\yi{简二十七}{地热衰减，二百年间下降百分之三点一。开头慢，最后急。一急，气先断，水其次，粟再次。最后一年，气路断了三次，又接上三次，吏把这事报上去。主事的人说：再断一次，就清点人口。}

\yi{简二十八}{外工三百二十人，不计入装船的人数。留下。留下的，给三个月的粮。}

\yi{简二十九}{六批的数目：第一批六千，第二批八千，第三批一万，第四批一万三千，第五批一万八千，第六批六万五千。共十二万。}

\yi{简三十}{六批不是一时的事。前五批在年度计划之内。第六批在年度计划之外。计划之内的有次序；计划之外的没有次序。}

\yi{简三十一}{关门那天，没有人送，没有哭声。关门的七个人，回头看了一眼。路上没有脚印，有痕迹的是风。}

\yi{简三十二}{之后三百七十二年，有人来过。门打不开，在墙上题了名字就走了。题名的三个人，其中一个的名字看不清。}

\yi{简三十三}{这不是月亮，这是一条船。船已经不能走了。我们守它，守了三代。守的人没有俸禄，也没有人接替。}

\yi{简三十四}{大明万历十七年，内官监奉旨采办，到这里就返回了。同行二百七十人，回来的有一百四十一人。}

\yi{简三十五}{缺。有十二个，不能穷尽。穷尽的有三个，剩下九个丢了。十二是三的四倍。三个穷尽了，九个丢了，丢的不留。}

\yi{简三十六}{一件东西，圆形，直径不到二寸，重若干。不是金属，不是石头。敲它不响。来源不明，暂时不归类。}

\yi{简三十七}{字到这里为止。此后没有字。}
\endgroup
